%% ============================================
%% ================ Preambule =================
%% ============================================
\documentclass[]{scrartcl}  
\usepackage[margin = 0.5in]{geometry}

\usepackage[pdftex,unicode,
colorlinks=true,
linkcolor=blue]{hyperref}
\usepackage[T1,T2A]{fontenc}
\usepackage[utf8]{inputenc}
\usepackage[english,russian]{babel}
\usepackage{amssymb,amsmath}

%% ============================================
%% ================ Info ======================
%% ============================================
\title{Investigating Gradient Inconsistency in Low-Rank Adaptation of Vision Models (LoRA-Pro)}
\author{
\begin{tabular}{c c}
\shortstack{Готовский\\Геннадий} & \shortstack{Габидулин\\Андрей} \\
\texttt{gotovskii.gd@phystech.edu} & \texttt{gabidulin.aa@phystech.edu} \\[0.75em]
\multicolumn{2}{c}{\shortstack{Сивачев\\Леонид}} \\
\multicolumn{2}{c}{\texttt{sivachev.ld@phystech.edu}}
\end{tabular}
}
\date{\today}

\begin{document}


\maketitle

\begin{abstract}
Проект посвящён сравнению методов параметр-эффективного дообучения Vision Transformer / CLIP-ViT для задач классификации изображений. Основной фокус --- проверка гипотезы LoRA-Pro о том, что часть разрыва между LoRA и full fine-tuning связана не только с низким рангом обновления, но и с несовпадением эффективного low-rank градиента с полным градиентом fine-tuning. В качестве базовых методов будут рассмотрены zero-shot/frozen baseline, linear probe, full fine-tuning, LoRA, DoRA и LoRA-Pro. Эксперименты планируется проводить на CLIP-ViT-B/16 и датасетах CIFAR-100, EuroSAT и DTD; Food-101 будет рассматриваться как расширение при наличии вычислительного бюджета. Итогом проекта станут воспроизводимый код, таблицы и графики сравнения методов по accuracy, gap к full fine-tuning, числу обучаемых параметров, памяти и времени обучения.
\end{abstract}

\section{Идея}

В рамках данного проекта планируется исследовать методы параметр-эффективного дообучения (\textit{parameter-efficient fine-tuning}, PEFT) для моделей компьютерного зрения, в частности для архитектуры \textit{Vision Transformer} (ViT). Основной интерес представляет метод \textbf{LoRA-Pro}, в котором рассматривается проблема \textit{gradient mismatch}: стандартное обновление низкоранговых матриц в методе LoRA не эквивалентно шагу градиентного спуска в исходном полном пространстве параметров.

Идея проекта состоит в том, чтобы проверить, позволяет ли корректировка low-rank updates, предложенная в LoRA-Pro, улучшить качество parameter-efficient fine-tuning в задачах классификации изображений. Для этого планируется реализовать и сравнить несколько подходов к адаптации одной и той же предобученной vision-модели: \textbf{LoRA}, \textbf{DoRA} и \textbf{LoRA-Pro}. В качестве эталона будет рассмотрен режим \textbf{Full Fine-Tuning}.

Экспериментальная часть проекта будет построена на готовых датасетах классификации изображений: в качестве основных будут использоваться \texttt{CIFAR-100}, \texttt{EuroSAT} и \texttt{DTD}, а \texttt{Food-101} будет рассматриваться как дополнительный датасет при наличии вычислительного бюджета. В качестве основной backbone-модели планируется использовать \texttt{CLIP-ViT-B/16}, поскольку LoRA-Pro уже оценивался на image classification с CLIP-ViT-B/16, а сама модель является естественным примером vision foundation model. Основной вопрос проекта можно сформулировать следующим образом: \textbf{сокращает ли LoRA-Pro разрыв между parameter-efficient fine-tuning и полным fine-tuning в задачах компьютерного зрения по сравнению со стандартным LoRA и DoRA?}

\subsection{LoRA как базовый метод low-rank адаптации}

Одним из основных baseline-методов в проекте будет LoRA
(\textit{Low-Rank Adaptation})~\cite{hu2022lora}. Метод был предложен как способ сделать
дообучение больших Transformer-моделей существенно дешевле по числу обучаемых
параметров, памяти и стоимости хранения checkpoint-ов. Проблема полного
fine-tuning состоит в том, что для каждой downstream-задачи необходимо
обновлять и хранить все параметры модели. Для больших моделей это приводит к
большим затратам памяти: помимо самих весов необходимо хранить градиенты и
состояния оптимизатора, например первые и вторые моменты в Adam.

LoRA исходит из следующей гипотезы: при адаптации предобученной модели к
конкретной задаче не обязательно существенно менять всю матрицу весов. Полезное
изменение весов может лежать в низкоразмерном подпространстве, то есть иметь
малый эффективный ранг. Важно, что речь идет не о низком ранге самой
предобученной матрицы $W_0$, а именно о низком ранге изменения $\Delta W$.

Пусть
\[
    W_0 \in \mathbb{R}^{d \times k}
\]
--- предобученная матрица некоторого dense-слоя. В полном fine-tuning мы
разрешаем произвольное обновление:
\[
    W = W_0 + \Delta W,
    \qquad
    \Delta W \in \mathbb{R}^{d \times k}.
\]
LoRA ограничивает это обновление низкоранговой параметризацией:
\[
    \Delta W = BA,
\]
где
\[
    B \in \mathbb{R}^{d \times r},
    \qquad
    A \in \mathbb{R}^{r \times k},
    \qquad
    r \ll \min(d,k).
\]
Тогда
\[
    \operatorname{rank}(\Delta W)
    =
    \operatorname{rank}(BA)
    \leq r.
\]
Иными словами, вместо обучения всех $dk$ элементов матрицы $W$ обучаются только
\[
    r(d+k)
\]
параметров в матрицах $A$ и $B$.

С точки зрения forward pass стандартный линейный слой
\[
    h = W_0x
\]
заменяется на
\[
    h = W_0x + \frac{\alpha}{r}BAx,
\]
где $\alpha$ --- масштабирующий коэффициент. При этом $W_0$ остается
замороженной, а оптимизация проводится только по малым матрицам $A$ и $B$.
Произведение $BA$ не обязательно явно строить как большую матрицу: вычисление
можно выполнить последовательно через bottleneck-размерность $r$:
\[
    x \in \mathbb{R}^{k}
    \xrightarrow{\;A\;}
    Ax \in \mathbb{R}^{r}
    \xrightarrow{\;B\;}
    BAx \in \mathbb{R}^{d}.
\]
Это объясняет, почему LoRA воздействует на большую матрицу весов, но обучает
малое число независимых параметров.

Мотивация low-rank подхода связана с SVD-разложением. Если бы после полного
fine-tuning мы получили некоторое изменение $\Delta W$, то его можно было бы
представить в виде
\[
    \Delta W
    =
    \sum_{i=1}^{\operatorname{rank}(\Delta W)}
    \sigma_i u_i v_i^\top.
\]
Если основные изменения сосредоточены в нескольких первых сингулярных
направлениях, то $\Delta W$ можно хорошо аппроксимировать матрицей малого
ранга. LoRA не вычисляет SVD напрямую и не строит лучшее SVD-приближение.
Вместо этого метод сразу ищет обновление в классе низкоранговых матриц
$\Delta W = BA$, оптимизируя $A$ и $B$ под downstream loss.

Авторы LoRA показали, что на ряде NLP-задач такой подход позволяет получить
качество, сравнимое с full fine-tuning, при гораздо меньшем числе обучаемых
параметров. Для GPT-3 175B они сообщают уменьшение числа обучаемых параметров
на порядки, а также отсутствие дополнительной задержки на inference, поскольку
после обучения LoRA-добавку можно слить с исходной матрицей:
\[
    W_{\text{merged}}
    =
    W_0 + \frac{\alpha}{r}BA.
\]
После такого слияния модель выполняет обычное матричное умножение, как и при
полном fine-tuning.

\subsection{Почему LoRA релевантна для Vision Transformer}

Хотя LoRA изначально была подробно исследована для языковых Transformer-моделей,
сама конструкция применима к любым dense-слоям. Vision Transformer также состоит
из Transformer-блоков, где изображение представляется как последовательность
patch-токенов, а основные вычисления выполняются в self-attention и MLP-модулях.
В каждом attention-блоке есть матрицы проекций
\[
    W_q,\quad W_k,\quad W_v,\quad W_o.
\]
К ним можно добавить low-rank обновления:
\[
    W_q = W_{q,0} + B_qA_q,
    \qquad
    W_v = W_{v,0} + B_vA_v,
\]
и аналогично для $W_k$ и $W_o$.

В оригинальной работе LoRA авторы отдельно исследовали, к каким attention-матрицам
лучше применять low-rank адаптацию. Их эксперименты показали, что адаптация
$W_q$ и $W_v$ часто оказывается эффективной: $W_q$ отвечает за то, какие запросы
делает токен к остальному контексту, а $W_v$ --- за то, какая информация передается
через attention. Поэтому в нашем проекте естественным первым вариантом будет
LoRA на $W_q$ и $W_v$. Однако для Vision Transformer это не является заранее
доказанным фактом: изображения, patch-токены и vision-задачи отличаются от
языковых последовательностей. Поэтому мы будем отдельно проверять, какие матрицы
ViT выгоднее адаптировать.

Планируется сравнить следующие варианты:
\begin{enumerate}
    \item основной режим: LoRA/DoRA/LoRA-Pro на всех linear-модулях visual backbone, кроме embedding, normalization и classification head;
    \item ablation-режим: адаптация только attention-проекций $W_q,W_v$;
    \item ablation-режим: адаптация всех attention-проекций $W_q,W_k,W_v,W_o$;
    \item основной ранг $r=8$ и абляция по $r \in \{4,8,16,32\}$ при наличии вычислительного бюджета.
\end{enumerate}
\section{Литературный обзор и место проекта}

В исходной работе по Vision Transformer показано, что изображение можно представить как последовательность patch-токенов и применять к ней обычный Transformer encoder~\cite{dosovitskiy2021vit}. Это делает ViT естественной архитектурой для проверки PEFT-методов, изначально разработанных для языковых Transformer-моделей. CLIP дополнительно показывает, что visual encoder, обученный на image-text парах, может переноситься на новые классы через текстовые prompts~\cite{radford2021clip}; поэтому CLIP-ViT является удобным backbone для экспериментов с downstream image classification.

LoRA является базовым методом low-rank адаптации: она замораживает предобученные веса и обучает только низкоранговую добавку $BA$~\cite{hu2022lora}. Это резко уменьшает число обучаемых параметров и позволяет после обучения слить добавку с исходным весом, не добавляя latency на inference. Однако LoRA обычно уступает full fine-tuning, особенно если low-rank ограничение или оптимизационная траектория оказываются слишком жёсткими.

DoRA предлагает альтернативное объяснение части gap-а между LoRA и full fine-tuning: стандартная LoRA одним обновлением $BA$ одновременно меняет magnitude и direction весов~\cite{liu2024dora}. Поэтому DoRA раскладывает вес на масштаб и направление, обучая magnitude отдельно, а direction --- через low-rank update. В отличие от DoRA, LoRA-Pro фокусируется не на параметризации веса, а на оптимизации: стандартный шаг LoRA можно интерпретировать как full fine-tuning с низкоранговым equivalent gradient~\cite{wang2025lorapro}. LoRA-Pro корректирует градиенты матриц $A$ и $B$, чтобы этот equivalent gradient лучше приближал полный градиент.

AdaLoRA занимает соседнее место в ландшафте: она адаптивно распределяет параметрический бюджет между слоями и фактически отвечает на вопрос ``где и какой ранг нужен''~\cite{zhang2023adalora}. Наш проект отвечает на другой вопрос: при фиксированном или контролируемом rank budget важнее ли для ViT-классификации стандартная low-rank параметризация, decomposition magnitude/direction или коррекция градиента.

Таким образом, проект находится на пересечении трёх линий работ: ViT/CLIP как vision foundation models, LoRA-подобные PEFT-методы и новые попытки объяснить и уменьшить gap между PEFT и full fine-tuning. Основной исследовательский пробел, который мы проверяем экспериментально, состоит в том, насколько идеи DoRA и LoRA-Pro переносятся на чистые задачи image classification с ViT/CLIP-ViT.

Из этого обзора следует экспериментальная постановка проекта: сравнивать нужно не только LoRA и full fine-tuning, но и методы, которые объясняют gap между ними по-разному. DoRA проверяет гипотезу о роли разложения весов на magnitude и direction, а LoRA-Pro проверяет гипотезу о роли gradient mismatch. Поэтому отрицательный результат также будет содержательным: если LoRA-Pro или DoRA не дают выигрыша на ViT/CLIP-ViT, это укажет, что в выбранной vision-постановке основным ограничением может быть сам low-rank budget, выбор target modules или датасетный сдвиг.

\subsection{Роль AdaLoRA и почему она не является основным методом}

В литературе существуют и более сложные low-rank методы, например AdaLoRA,
который адаптивно распределяет параметрический бюджет между слоями и может
динамически менять эффективный ранг разных адаптеров. Однако в данном проекте
основной акцент делается не на адаптивном выборе ранга, а на сравнении разных
способов оптимизации и параметризации low-rank адаптеров: LoRA, DoRA и LoRA-Pro.

Поэтому AdaLoRA пока не включается в основной набор методов, но будет рассмотрена
в литературном обзоре как важный related work. Иначе проект становится слишком
широким: пришлось бы одновременно исследовать адаптивный выбор ранга, выбор
слоев, разложение весов на magnitude/direction и проблему gradient mismatch.

В рамках текущей постановки LoRA используется как базовый low-rank baseline,
DoRA --- как модификация, учитывающая разложение веса на направление и масштаб,
а LoRA-Pro --- как метод, специально мотивированный проблемой некорректной
оптимизации low-rank факторов. AdaLoRA может быть добавлен как дополнительный
baseline при наличии времени и вычислительных ресурсов.

% --- DoRA subsection above ---
\subsection{DoRA: разложение веса на magnitude и direction}

Следующим методом для сравнения будет DoRA (\textit{Weight-Decomposed Low-Rank
Adaptation})~\cite{liu2024dora}. Его мотивация возникает из ограничения стандартной LoRA. В LoRA
обновление
\[
    W' = W_0 + BA
\]
одновременно меняет и масштаб весов, и их направление. Поэтому одна и та же
low-rank добавка должна отвечать сразу за два разных эффекта: насколько сильно
изменить весовой вектор и в какую сторону его повернуть. В работе DoRA
показывается, что такое совместное обновление может быть менее гибким, чем full
fine-tuning.

Идея DoRA состоит в том, чтобы явно разложить каждый столбец матрицы весов на
длину и направление. Пусть
\[
    W = [w_1,\ldots,w_k] \in \mathbb{R}^{d \times k}.
\]
Тогда каждый столбец можно записать как
\[
    w_j = \|w_j\| \frac{w_j}{\|w_j\|}.
\]
В матричной форме это записывается как
\[
    W = m \frac{V}{\|V\|_c},
\]
где $m \in \mathbb{R}^{1 \times k}$ --- вектор масштабов столбцов, $V \in
\mathbb{R}^{d \times k}$ --- матрица направлений, а $\|V\|_c$ обозначает
постолбцовую норму. Это не обычное матричное умножение, а scaling по столбцам:
каждый нормированный столбец $v_j / \|v_j\|$ умножается на свой коэффициент
$m_j$.

В DoRA предобученный вес инициализируется как
\[
    m = \|W_0\|_c, \qquad V = W_0.
\]
После этого масштаб $m$ обучается напрямую, а направление обновляется с помощью
LoRA-style low-rank добавки:
\[
    V' = W_0 + BA.
\]
Итоговый вес имеет вид
\[
    W'
    =
    m \frac{W_0 + BA}{\|W_0 + BA\|_c}.
\]
Если в начале обучения $B=0$, то $BA=0$, и поэтому
\[
    W' = \|W_0\|_c \frac{W_0}{\|W_0\|_c} = W_0.
\]
Значит DoRA, как и LoRA, стартует ровно с предобученной модели и не разрушает
исходные веса на первом шаге.

Такое разложение важно потому, что полное дообучение умеет независимо менять
длину и направление весовых векторов. В терминах DoRA это означает, что full
fine-tuning может делать большой magnitude update при малом изменении direction
или, наоборот, сильно менять direction почти без изменения magnitude. Стандартная
LoRA чаще связывает эти два эффекта, поскольку один и тот же $BA$ отвечает и за
масштаб, и за направление. DoRA разделяет эти роли: параметр $m$ отвечает за
magnitude, а low-rank добавка $BA$ отвечает за directional update.

Для проекта это важно по двум причинам. Во-первых, DoRA сохраняет главное
практическое преимущество LoRA: после обучения вес можно заранее собрать в
обычную матрицу $W'$, поэтому дополнительной задержки на inference не возникает.
Во-вторых, DoRA проверяет более тонкую гипотезу, чем обычная LoRA: возможно, для
адаптации ViT важен не только малый ранг изменения $\Delta W$, но и отдельный
контроль масштаба и направления весов. Поэтому в экспериментах мы будем
сравнивать LoRA и DoRA при одинаковом backbone, одинаковых датасетах и близком
числе обучаемых параметров.

В рамках проекта планируется проверить следующие варианты DoRA:
\begin{enumerate}
    \item DoRA на тех же attention-проекциях, что и LoRA, прежде всего $W_q$ и $W_v$;
    \item DoRA на всех attention-проекциях $W_q,W_k,W_v,W_o$;
    \item сравнение DoRA и LoRA при одинаковых значениях ранга $r$;
    \item сравнение DoRA и LoRA при близком числе обучаемых параметров.
\end{enumerate}

Если DoRA будет стабильно превосходить LoRA на ViT-классификации при близком
parameter budget, это будет означать, что разделение magnitude и direction
действительно полезно не только для языковых и vision-language моделей, но и для
чистых задач компьютерного зрения. Если же выигрыша не будет, это тоже будет
содержательным результатом: возможно, для image classification на ViT основным
ограничением является не связка magnitude/direction, а выбор слоев, ранг $r$ или
сам характер low-rank ограничения.


% --- LoRA-Pro subsection inserted here ---
\subsection{LoRA-Pro: коррекция градиента в low-rank адаптации}

Третьим методом для сравнения будет LoRA-Pro~\cite{wang2025lorapro}. В отличие от DoRA, который
модифицирует параметризацию весов через разложение на magnitude и direction,
LoRA-Pro рассматривает другую проблему стандартной LoRA: не только какие
обновления весов можно представить в low-rank форме, но и каким градиентом
метод фактически движется в процессе оптимизации.

В стандартной LoRA вес адаптируемого слоя записывается как
\[
    W = W_0 + sBA,
\]
где $s$ --- масштабирующий коэффициент. Если бы мы проводили full fine-tuning,
то обновление матрицы $W$ выполнялось бы по полному градиенту
\[
    g = \frac{\partial \mathcal{L}}{\partial W}.
\]
Однако в LoRA матрица $W$ не является напрямую обучаемым параметром: обучаются
только матрицы $A$ и $B$. Поэтому малое изменение веса имеет вид
\[
    dW = s\,dB\,A + s\,B\,dA.
\]
Если обозначить градиенты, по которым обновляются low-rank матрицы, через
$g_A$ и $g_B$, то соответствующий виртуальный градиент в пространстве полной
матрицы $W$ можно записать как
\[
    \widetilde{g} = sBg_A + sg_BA.
\]
LoRA-Pro называет эту величину \textit{equivalent gradient}. Она показывает,
что оптимизация LoRA эквивалентна full fine-tuning, но не с полным градиентом
$g$, а с некоторым low-rank градиентом $\widetilde{g}$.

Отсюда возникает проблема gradient mismatch:
\[
    \widetilde{g} \neq g.
\]
Даже если low-rank параметризация $BA$ достаточно выразительна для представления
нужного изменения весов, стандартная LoRA может двигаться к этому изменению по
неудачной траектории. Иными словами, LoRA ограничивает не только пространство
допустимых обновлений, но и эффективное направление оптимизационного шага.

Авторы LoRA-Pro предлагают не менять структуру адаптера, а корректировать градиенты
матриц $A$ и $B$ так, чтобы индуцированный equivalent gradient был как можно
ближе к градиенту full fine-tuning:
\[
    \min_{g_A,g_B}
    \left\|sBg_A + sg_BA - g\right\|_F^2.
\]
Интуитивно это можно понимать как выбор наилучшего доступного шага внутри
текущего low-rank tangent-пространства. LoRA-Pro не выводит модель за пределы
low-rank ограничения: обновление по-прежнему реализуется через матрицы $A$ и
$B$. Однако среди всех шагов, которые можно выразить через текущие $A$ и $B$,
метод выбирает тот, который лучше всего приближает шаг полного fine-tuning.

Это отличие важно для нашего проекта. LoRA и DoRA проверяют, насколько хорошо
low-rank параметризация весов подходит для адаптации ViT. LoRA-Pro проверяет
другую гипотезу: возможно, часть разрыва между PEFT и full fine-tuning связана
не только с ограниченной выразительностью low-rank обновления, но и с тем, что
стандартные градиенты по $A$ и $B$ задают недостаточно хорошее направление
движения в пространстве полной матрицы.

В экспериментах LoRA-Pro будет рассматриваться как метод, направленный именно
на проблему оптимизации low-rank адаптеров. Планируется сравнить его со
стандартной LoRA при одинаковом ранге $r$ и одинаковом выборе адаптируемых
матриц. Если LoRA-Pro будет превосходить LoRA при близком числе обучаемых
параметров, это будет свидетельствовать о том, что gradient mismatch является
существенным фактором для ViT-классификации. Если же выигрыша не будет, это
может означать, что в выбранной постановке основным bottleneck является не
оптимизационная траектория, а само low-rank ограничение, выбор слоев или
датасетный сдвиг.

\subsection{Исследовательская гипотеза}

\begin{quote}
Для адаптации предобученного Vision Transformer к downstream image classification
задачам может быть достаточно низкоранговых обновлений отдельных dense-матриц,
но эффективность такой адаптации зависит от выбора слоев, ранга $r$ и способа
оптимизации low-rank параметров.
\end{quote}

LoRA в этом проекте играет роль базового метода, относительно которого будут
оцениваться более современные подходы DoRA и LoRA-Pro. Это позволит отделить
два вопроса: достаточно ли самой идеи low-rank адаптации для ViT, и дает ли более
аккуратная параметризация или оптимизация low-rank обновлений дополнительный
выигрыш относительно стандартной LoRA.

\section{Постановка задачи}

Пусть дана предобученная vision-модель $f_{\Phi_0}$ с backbone на основе ViT/CLIP-ViT и downstream-датасет классификации изображений
\[
    \mathcal{D}=\{(x_i,y_i)\}_{i=1}^{N}.
\]
Цель проекта --- сравнить несколько способов адаптации одной и той же предобученной модели к новой задаче при контролируемом числе обучаемых параметров. Для всех методов используется один и тот же classification loss
\[
    \mathcal{L}(\Phi)
    =
    \frac{1}{N}
    \sum_{i=1}^{N}
    \ell(f_{\Phi}(x_i),y_i),
\]
где $\ell$ --- cross-entropy loss для задачи классификации.

В режиме full fine-tuning оптимизация проводится по всем параметрам модели:
\[
    \min_{\Phi}\mathcal{L}(\Phi).
\]
Этот режим является сильным, но дорогим baseline: он задаёт верхнюю точку сравнения по качеству, но требует обучать и хранить все параметры backbone.

В режиме LoRA предобученные веса $\Phi_0$ замораживаются, а для выбранного множества матриц $\mathcal{S}$ обучаются только low-rank параметры
\[
    \Theta_{\mathrm{LoRA}}=\{A_l,B_l\}_{l\in\mathcal{S}}.
\]
Для каждой адаптируемой матрицы
\[
    W_l'
    =
    W_{l,0}+\frac{\alpha}{r}B_lA_l,
\]
поэтому задача имеет вид
\[
    \min_{\Theta_{\mathrm{LoRA}}}
    \mathcal{L}\bigl(\Phi_0+\Delta\Phi(\Theta_{\mathrm{LoRA}})\bigr).
\]
LoRA проверяет базовую гипотезу: достаточно ли низкорангового изменения весов для адаптации ViT/CLIP-ViT к новой image classification задаче.

В режиме DoRA структура low-rank адаптера дополняется разложением веса на magnitude и direction. Для каждой адаптируемой матрицы обучаются
\[
    \Theta_{\mathrm{DoRA}}=\{A_l,B_l,m_l\}_{l\in\mathcal{S}},
\]
а сам вес задаётся как
\[
    W_l'
    =
    m_l\frac{W_{l,0}+B_lA_l}{\|W_{l,0}+B_lA_l\|_c}.
\]
DoRA проверяет вторую гипотезу: возможно, часть gap-а между LoRA и full fine-tuning связана не только с low-rank ограничением, но и с тем, что стандартная LoRA одним обновлением одновременно меняет масштаб и направление весов.

\paragraph{Core optimization problem из LoRA-Pro.}
Ключевая оптимизационная задача, предложенная авторами LoRA-Pro, состоит в том, чтобы подобрать скорректированные градиенты low-rank факторов так, чтобы эквивалентный low-rank градиент был как можно ближе к full fine-tuning gradient:
\begin{equation}
    \min_{g_{A_l},g_{B_l}}
    \left\|sB_lg_{A_l}+sg_{B_l}A_l-g\right\|_F^2.
    \label{eq:lorapro_core}
\end{equation}
В нашем проекте эта задача не заявляется как новый метод; она используется как центральная идея статьи LoRA-Pro, которую мы хотим воспроизвести и проверить на ViT/CLIP-ViT классификации.

В режиме LoRA-Pro структура адаптера остаётся LoRA-подобной, но изменяется оптимизационный шаг. Пусть
\[
    g=\frac{\partial \mathcal{L}}{\partial W_l}
\]
--- полный градиент по адаптируемой матрице, а
\[
    \widetilde{g}=sB_lg_{A_l}+sg_{B_l}A_l
\]
--- equivalent gradient, индуцированный градиентами low-rank матриц. LoRA-Pro корректирует $g_{A_l}$ и $g_{B_l}$ так, чтобы уменьшить расхождение $\|\widetilde{g}-g\|_F$, как в задаче~\eqref{eq:lorapro_core}.
Тем самым LoRA-Pro проверяет третью гипотезу: возможно, часть потери качества стандартной LoRA связана не только с ограниченным классом low-rank обновлений, но и с тем, что стандартные градиенты по $A_l$ и $B_l$ задают неудачное направление движения внутри допустимого low-rank tangent-пространства.

Итак, задача проекта состоит в экспериментальном сравнении четырёх режимов:
\[
    \text{Full Fine-Tuning},\quad
    \text{LoRA},\quad
    \text{DoRA},\quad
    \text{LoRA-Pro}
\]
на одинаковых backbone-моделях, датасетах и target modules. Основной измеряемый вопрос: какой фактор сильнее влияет на gap к full fine-tuning в ViT/CLIP-ViT классификации --- само low-rank ограничение, разложение magnitude/direction или gradient mismatch.

% --- Experiments section ---
\section{Планируемые эксперименты}

Основной backbone --- \texttt{CLIP-ViT-B/16} из \texttt{HuggingFace Transformers}, загружаемый как checkpoint \texttt{openai/clip-vit-base-patch16} из Hugging Face Hub~\cite{huggingfaceclipvit}. Дополнительно, при наличии времени, будет рассмотрен supervised ViT-B/16 из \texttt{timm}. Основные датасеты: \texttt{CIFAR-100}, \texttt{EuroSAT} и \texttt{DTD}; \texttt{Food-101} и датасеты из LoRA-Pro image-classification benchmark будут использоваться только как расширение.

\subsection{Источники кода, моделей и данных}

Для практической реализации проекта планируется опираться не только на статьи, но и на открытые реализации и модельные карточки. Основной backbone будет загружаться из HuggingFace Hub:
\[
    \texttt{openai/clip-vit-base-patch16}.
\]
Эта модельная карточка фиксирует используемую архитектуру CLIP-ViT-B/16 и позволяет воспроизводимо загрузить image encoder и text encoder для zero-shot baseline и дальнейшего fine-tuning~\cite{huggingfaceclipvit}.

В качестве базовой инфраструктуры для LoRA и PEFT будет использоваться библиотека HuggingFace PEFT и её image-classification examples~\cite{huggingfacepeftlora}.
Это позволит быстро получить корректный LoRA baseline, проверить число обучаемых параметров и затем заменить стандартное обновление LoRA на LoRA-Pro.

Для LoRA-Pro будет изучена открытая реализация авторов~\cite{loraprogithub}.
Её планируется использовать как основной reference implementation для gradient correction. Если код окажется жёстко привязан к языковым моделям или конкретным HuggingFace modules, будет реализована минимальная версия LoRA-Pro для linear-слоёв CLIP-ViT с сохранением основной формулы equivalent gradient.

Для DoRA будет использоваться официальный репозиторий NVIDIA~\cite{doragithub}.
Он нужен не только для возможного запуска DoRA baseline, но и для проверки деталей параметризации magnitude/direction и практического merge весов перед inference.

Датасеты \texttt{CIFAR-100}, \texttt{EuroSAT} и \texttt{DTD} будут загружаться через \texttt{torchvision} или HuggingFace Datasets в зависимости от удобства preprocessing pipeline. Все конфигурации запуска, версии библиотек и ссылки на используемые checkpoint-и будут сохранены в репозитории проекта, чтобы результаты можно было воспроизвести.

После фазы прототипирования основной pipeline будет расширен следующим образом:
\begin{enumerate}
    \item загрузка \texttt{CLIP-ViT-B/16} и датасетов \texttt{CIFAR-100}, \texttt{EuroSAT}, \texttt{DTD};
    \item реализация zero-shot CLIP baseline через prompts вида \texttt{a photo of a \{class\}};
    \item запуск linear probe / frozen visual encoder baseline;
    \item запуск full fine-tuning visual backbone;
    \item запуск LoRA при $r=8$, $\alpha=16$;
    \item запуск DoRA при том же rank и тех же target modules;
    \item запуск LoRA-Pro при той же LoRA-конфигурации;
    \item сравнение качества, числа обучаемых параметров, времени обучения и peak GPU memory.
\end{enumerate}


После реализации LoRA-базовой линии планируется добавить DoRA и сравнить его с LoRA в тех же
условиях. Затем будет добавлен LoRA-Pro как метод, специально направленный на проблему
optimization/gradient mismatch в low-rank адаптерах.

\subsection{План работ}

Проект рассчитан примерно на три недели активной экспериментальной работы. На первом этапе будет собран минимальный pipeline на subset \texttt{CIFAR-100}: zero-shot CLIP, linear probe, LoRA при $r=8$ и проверка корректного подсчёта trainable parameters. На втором этапе будут запущены основные эксперименты на \texttt{CIFAR-100} и \texttt{EuroSAT}: full fine-tuning, LoRA, DoRA и LoRA-Pro при одинаковом backbone, target modules и rank. На третьем этапе будут проведены абляции по rank $r \in \{4,8,16,32\}$ и по множеству адаптируемых модулей: $W_q,W_v$, все attention-проекции и все linear-модули visual backbone. Если останется вычислительный бюджет, будет добавлен \texttt{DTD} и/или \texttt{Food-101}.

% --- Prototyping phase ---
\subsection{Фаза прототипирования}

На этапе прототипирования был собран базовый пайплайн на \texttt{HuggingFace Transformers}, \texttt{torchvision} и Hugging Face Hub: загрузка checkpoint-а \texttt{openai/clip-vit-base-patch16}~\cite{huggingfaceclipvit}, подготовка \texttt{CIFAR-100}, генерация текстовых prompts вида \texttt{a photo of a \{class\}} и zero-shot оценка без обучения backbone. В этой конфигурации zero-shot accuracy на \texttt{CIFAR-100} составила $62.8\%$, что мы используем как первую воспроизводимую точку отсчёта для последующих PEFT-экспериментов.

Также был реализован базовый класс LoRA-адаптера для linear-слоёв visual encoder и проверен подсчёт trainable parameters для ранга $r=8$. Основная техническая сложность на этом этапе оказалась связана с тем, что стандартный \texttt{Trainer} в \texttt{HuggingFace} неудобен для диагностики LoRA-Pro: нужно явно перехватывать градиенты по low-rank факторам и сравнивать их с full-gradient proxy по исходной матрице. Поэтому для прототипа был подготовлен кастомный training loop, в котором forward pass, backward pass и сбор diagnostic-метрик контролируются вручную.

Прототип кода для вычисления equivalent gradient
\[
    \widetilde{g}=sB_lg_{A_l}+sg_{B_l}A_l
\]
прошёл sanity-check на малых случайных матрицах: проверялись размерности, корректность Frobenius-нормы $\|\widetilde{g}-g\|_F$ и совпадение численного масштаба при изменении коэффициента $s$. Предварительные замеры показали, что LoRA при $r=8$ существенно сокращает число обучаемых параметров относительно full fine-tuning; следующий шаг --- запустить тот же прототип на \texttt{EuroSAT} и добавить корректировку градиента LoRA-Pro.

% --- Outcomes section ---
\section{Outcomes}

Конкретными выходами проекта будут:
\begin{enumerate}
    \item воспроизводимый код обучения и оценки CLIP-ViT-B/16 с full fine-tuning, LoRA, DoRA и LoRA-Pro;
    \item конфигурационные файлы для всех основных запусков;
    \item список использованных открытых источников кода и checkpoint-ов: HuggingFace CLIP-ViT-B/16~\cite{huggingfaceclipvit}, HuggingFace PEFT~\cite{huggingfacepeftlora}, LoRA-Pro GitHub~\cite{loraprogithub} и DoRA GitHub~\cite{doragithub};
    \item таблицы результатов по датасетам, методам, rank и target modules;
    \item графики train/validation loss, accuracy vs trainable parameters, gap-to-full-fine-tuning и memory/time comparison;
    \item короткий аналитический отчёт о том, какой фактор оказался важнее для ViT-классификации: low-rank parameterization, magnitude/direction decomposition или gradient correction;
    \item при наличии времени --- диагностический эксперимент по величине $\|\widetilde{g}-g\|_F$ на небольших probe-batches для LoRA и LoRA-Pro.
\end{enumerate}
\section{Метрики качества}

Для объективного сравнения методов будут использоваться следующие метрики:
\begin{enumerate}
    \item validation accuracy;
    \item число обучаемых параметров;
    \item доля обучаемых параметров относительно полной модели;
    \item скорость сходимости по эпохам;
    \item максимальное потребление GPU-памяти;
    \item gap относительно full fine-tuning:
    \[
        \Delta_{\text{gap}}
        =
        \operatorname{Acc}_{\text{Full FT}}
        -
        \operatorname{Acc}_{\text{PEFT}}.
    \]
\end{enumerate}

Дополнительно для LoRA-Pro, если вычислительно возможно, будет измеряться diagnostic-метрика
\[
    E_{\mathrm{grad}}
    =
    \|\widetilde{g} - g\|_F,
\]
где $g$ --- полный градиент по выбранной матрице на probe-batch, а $\widetilde{g}$ --- equivalent gradient, индуцированный LoRA/LoRA-Pro. Эта метрика не является основной, но позволяет проверить саму мотивацию LoRA-Pro.

Отрицательный результат также будет содержательным: если LoRA окажется существенно хуже full
fine-tuning на выбранных vision-задачах, это будет означать, что low-rank ограничение недостаточно
выразительно в данной постановке или требует более аккуратного выбора слоев и ранга.



% --- Bibliography section ---
\renewcommand{\refname}{Список литературы}
\bibliographystyle{plain}
\bibliography{biblio}

\end{document}
